\clearpage
\setcounter{page}{1}
\maketitlesupplementary

% \part*{Supplementary Material}
% %\begin{center}
% %	{\LARGE \textbf{Appendix}}
% %\end{center}
% {
% % \setlength{\parskip}{-0em}
% \startcontents[sections]
% \printcontents[sections]{ }{1}{}
% }
% \clearpage

\section{Example of Curated Data}
\cref{tab:cot-example} illustrates how we guide the model to inject physical reasoning: starting with a concise scenario description, we enforce the relevant law, clarify the expected visual outcome, and provide a step‐by‐step rationale that links microgravity and surface tension to the observed behavior of the liquid. This structured example demonstrates the depth and clarity of our CoT annotations.
\begin{table*}[h]
  \caption{Example from our curated physics‐focused Chain‐of‐Thought dataset. The original caption sets the scene in microgravity, the annotated law (‘Lack of gravity’) defines the physical constraint, the explicit caption specifies the outcome (spherical droplets floating), and the chain of thought.}
  \centering
  \begin{tabular}{@{}p{3cm}p{10cm}@{}}
    \toprule
    \textbf{Caption} & A can of soda is slowly poured out in the space station, releasing the liquid into the surrounding area. \\
    \midrule
    \textbf{Physical Law} & Lack of gravity. \\
    \midrule
    \textbf{Explicit Caption} & A can of soda is slowly poured out in the space station, releasing the liquid into the surrounding area. The liquid forms spherical droplets and floats freely in the air. \\
    \midrule
    \textbf{Chain of Thought} & In microgravity (or “zero-g” conditions), there is no dominant downward force pulling the liquid. Instead of flowing downward as it would on Earth, the soda drifts away once it exits the can. Surface tension becomes the main force acting on the soda droplets, pulling the liquid into the smallest possible shape, a sphere. With no significant gravitational force to break them apart, these droplets remain spherical and continue floating until they contact another surface or droplet. Hence, in the absence of gravity, soda poured in a spacecraft forms and floats as spherical droplets. \\
    \bottomrule
  \end{tabular}
  \label{tab:cot-example}
\end{table*}


\section{Video Generator Settings}
The evaluated teams converged on a similar fast-preview configuration, typically employing around 50 diffusion inference steps with classifier-free guidance to balance efficiency and visual quality. Specifically, CogVideoX-5B/2B utilize the DPM-Solver scheduler (CogVideoXDPMScheduler) configured with classifier-free guidance set to 6, producing 24-frame video clips (6 seconds at 4 fps) at a resolution of 480 × 720 pixels, all executed in bfloat16 precision \cite{yang2024cogvideox}.
VideoCrafter2, by contrast, retains the classical DDIM sampler while increasing classifier-free guidance to 12, generating within the same 50-step inference budget video sequences of 24 frames (6 seconds at 4 fps) at 320 × 512 pixels in full fp32 precision \cite{wang2023lavie}.
Similarly, LaVie leverages its base UNet architecture with classifier-free guidance set to 12, matching VideoCrafter2’s inference settings—50 DDIM steps to produce 24-frame sequences (6 seconds at 4 fps) at 320 × 512 resolution in fp32 precision \cite{chen2024videocrafter2}.
All three diffusion-based models employ a cosine-based $\beta$ scheduling (Stable Diffusion’s standard schedule), utilize v-prediction parameterization, and explicitly follow a commonly adopted classifier-free guidance protocol with approximately 50 inference steps.

\section{Transferability of PhyPrompt to Different Video Generators}
\textbf{Results on Lavie.}
\cref{tab:enhanced_models} reports the performance of each prompt refinement strategy when generating videos with Lavie. Using the original prompts yields a baseline SA of 41.0\% and PC of 57.0\%, for a joint success rate of 29.2\%. Promptist underperforms the baseline, while simple instruction-tuned models (Qwen2.5) show modest gains in PC at the expense of SA. In contrast, PhyPrompt variants consistently improve both metrics – for example, PhyPrompt-7B raises PC to 63.2\% and joint success to 31.6\% – confirming its ability to inject physics-aware reasoning into prompts without degrading semantic alignment. GPT-4o and DeepSeek-V3 achieve near-baseline results, underscoring the unique advantage of our two-stage RL-trained rewriter on the Lavie generator.
\begin{table}[ht]
% \vspace{-2em}
  \centering
\caption{Semantic Adherence (SA), Physical Commonsense (PC), and joint SA\&PC scores for various prompt enhancement methods evaluated on the \textbf{Lavie} generator.}
  \label{tab:enhanced_models}
  \setlength{\tabcolsep}{3pt}
  \begin{tabular}{lccc}
    \toprule
    \textbf{Enhanced Model} & \textbf{SA (\%)} & \textbf{PC (\%)} & \textbf{PC \& SA (\%)} \\
    \midrule
    N/A (Original Prompt) & 41.0 & 57.0 & 29.2 \\
    Promptist             & 37.8 & 58.4 & 27.0 \\
    PhyT2V                & 41.2 & 60.8 & 30.0 \\
    \midrule
    Qwen2.5-1.5B          & 34.6 & 59.0 & 24.8 \\
    PhyPrompt-1.5B        & 40.2 & 61.2 & 30.6 \\
    Qwen2.5-3B            & 35.2 & 60.4 & 28.0 \\
    PhyPrompt-3B          & 41.4 & 60.8 & 30.2 \\
    Qwen2.5-7B            & 40.0 & 60.8 & 30.8 \\
    PhyPrompt-7B          & \textbf{41.8} & \textbf{63.2} & \textbf{31.6}\\
    \midrule
    GPT-4o                & 41.8 & 58.0 & 29.8 \\
    DeepSeek-V3           & 41.6 & 59.2 & 30.0 \\
    \bottomrule
  \end{tabular}
\end{table}

\textbf{Results on VideoCrafter2.}
\cref{tab:enhanced_vc2} reports results on VideoCrafter2. The original prompts achieve 45.6\% SA and 48.0\% PC for a 29.8\% joint score. Promptist underperforms the baseline, while Qwen2.5 variants trade off SA for higher PC (e.g., Qwen2.5-1.5B yields 35.0\% SA and 53.6\% PC). In contrast, PhyPrompt consistently boosts both metrics: the 1.5B variant raises joint success to 30.2\%, and the 7B variant attains the best scores (46.2\% SA, 66.4\% PC, 34.8\% joint). GPT-4o and DeepSeek-V3 deliver moderate PC improvements but fall short of PhyPrompt.
\begin{table}[ht]
  \centering
  % \vspace{-2em}
  \caption{Semantic Adherence (SA), Physical Commonsense (PC), and joint SA\&PC scores for various prompt enhancement methods evaluated on the \textbf{VideoCrafter2} generator.}
  \setlength{\tabcolsep}{3pt}
  \label{tab:enhanced_vc2}
  \begin{tabular}{lccc}
    \toprule
    \textbf{Enhanced Model} & \textbf{SA (\%)} & \textbf{PC (\%)} & \textbf{PC \& SA (\%)} \\
    \midrule
    N/A (Original Prompt) & 45.6 & 48.0 & 29.8 \\
    Promptist             & 43.8 & 51.4 & 27.2 \\
    PhyT2V                & 45.0 & 56.8 & 30.6 \\
    \midrule
    Qwen2.5-1.5B          & 35.0 & 53.6 & 27.8 \\
    PhyPrompt-1.5B        & 41.4 & 57.4 & 30.2 \\
    Qwen2.5-3B            & 40.0 & 55.8 & 28.0 \\
    PhyPrompt-3B          & 38.0 & 61.0 & 30.8 \\
    Qwen2.5-7B            & 45.8 & 58.0 & 31.2 \\
    PhyPrompt-7B          & \textbf{46.2} & \textbf{66.4} & \textbf{34.8} \\
    \midrule
    GPT-4o                & 38.4 & 60.0 & 30.0 \\
    DeepSeek-V3           & 40.0 & 60.4 & 32.8 \\
    \bottomrule
  \end{tabular}
\end{table}

\textbf{Results on CogVideoX-5B.}
\cref{tab:enhanced_cv5b} shows the performance of each prompt refinement strategy when driving the CogVideoX-5B generator. The original prompts achieve 50.2\% SA and 63.4\% PC (39.4\% joint). Promptist again underperforms the baseline, while Qwen2.5 models trade semantic fidelity for modest PC gains.
In contrast, PhyPrompt variants largely restore SA and push PC higher: PhyPrompt-7B matches the baseline SA at 50.2\% and attains the best PC of 69.2\%, yielding a 42.0\% joint score.
GPT-4o delivers the highest SA (52.6\%) but falls short of PhyPrompt on PC, and DeepSeek-V3 shows balanced but lower joint performance.
These results confirm PhyPrompt’s ability to inject physics-aware reasoning into prompts for CogVideoX-5B without degrading semantic.

\begin{table}[ht]
  \centering
  \caption{Semantic Adherence (SA), Physical Commonsense (PC), and joint SA\&PC scores for various prompt enhancement methods evaluated on \textbf{CogVideoX-5B}.}
  \setlength{\tabcolsep}{3pt}
  \label{tab:enhanced_cv5b}
  \begin{tabular}{lccc}
    \toprule
    \textbf{Enhanced Model}   & \textbf{SA (\%)} & \textbf{PC (\%)} & \textbf{PC\&SA (\%)} \\
    \midrule
    N/A (Original Prompt)     & 50.2             & 63.4             & 39.4                   \\
    Promptist                 & 46.4             & 63.0             & 37.4                   \\
    PhyT2V                    & 50.6             & 65.0             & 40.8                   \\
    \midrule
    Qwen2.5-1.5B              & 39.2             & 66.4             & 31.0                   \\
    PhyPrompt-1.5B            & 47.4             & 66.6             & 39.0                   \\
    Qwen2.5-3B                & 45.2             & 64.6             & 35.8                   \\
    PhyPrompt-3B              & 46.4             & 68.4             & 40.2                   \\
    Qwen2.5-7B                & 47.8             & 63.8             & 38.2                   \\
    PhyPrompt-7B              & 50.2             & \textbf{69.2}             & \textbf{42.0}                  \\
    \midrule
    GPT-4o                    & \textbf{52.6}             & 66.4             & 41.8                   \\
    DeepSeek-V3               & 50.6             & 63.6             & 40.8                   \\
    \bottomrule
  \end{tabular}
\end{table}


\section{Training Settings}
\textbf{SFT Training Setup.}
We initialize from Qwen2.5-Instruct (1.5B/3B/7B)~\cite{Qwen2024}, which has been pretrained on diverse corpora and instruction-tuned for dialogue.
To elicit explicit chain-of-thought reasoning, we insert token delimiters as follows:\texttt{<|im\_start|>think} and \texttt{<|im\_start|>answer}.
Each marker is surrounded by newlines to clearly separate “think” and “answer” stages.
We adopt the fine-tuning regimen of \cite{muennighoff2025simple}, optimizing the full model end-to-end for reasoning.
Training runs for $5$ epochs with a per-GPU batch size of $16$.
All computations use bfloat16 precision and a maximum sequence length of $4096$ tokens.
We employ AdamW \cite{loshchilov2019decoupled} with $\beta_{1}=0.9$, $\beta_{2}=0.95$, weight decay $1\times10^{-4}$, and gradient clipping at $1.0$.
The learning rate warms up linearly to $1\times10^{-5}$ over the first $5\%$ of steps and then decays to zero via a cosine schedule over the remaining steps.
We compute cross-entropy loss solely on reasoning traces and solutions, omitting question tokens. Fine-tuning is performed on a single node with $8\times\mathrm{A}100$ GPUs using Fully Sharded Data Parallel (FSDP) with CPU off-loading and gradient checkpointing; we save checkpoints every epoch and evaluate on a held-out validation set.


\textbf{GRPO Training Setup.}
We fine-tune the policy LLM for a total of 250 steps using a learning rate of $1\times10^{-6}$, decayed linearly to zero.
Training is conducted on a single node equipped with $8{\times}\mathrm{A}100$ GPUs, leveraging Fully-Sharded Data Parallel (FSDP) for efficient GPU memory usage, along with CPU off-loading and gradient checkpointing to further reduce memory footprint.
During policy roll-outs, we utilize vLLM to generate \textbf{4} candidate responses per prompt, employing temperature $t=1.0$ and top-$p$ sampling with $p=1.0$ to ensure diversity.
Input sequences are truncated to a maximum of 4096 tokens, while both generated responses and reasoning process are limited to at most 1024 tokens each.
Policy optimization is performed using Generalized Advantage Estimation (GAE) with parameters $\lambda=1$ and $\gamma=1$, alongside a KL-regularization term controlled by coefficient $\beta=0.01$ and a maximum gradient norm of 20.
We set the per-GPU batch size to 1, resulting in an overall batch size of 8 across all GPUs, and checkpoints are saved every 100 training steps.



\section{Model Architecture}
\cref{tab:details} contrasts the architectures and prompt capacities of five video-generation and evaluation models.
CogVideoX-5B and CogVideoX-2B \cite{wang2024cogvideodex} both pair a 3D causal VAE with a Transformer backbone (5B and 2B parameters, respectively), use a 4 096‐dimensional hidden state, and support long prompts of up to 226 tokens via RoPE or sinusoidal positional encodings.
By comparison, LaVie \cite{wang2023lavie} employs an 860 M-parameter cascaded latent diffusion pipeline, VideoCrafter2 \cite{chen2024videocrafter2} uses a 1.2B-parameter UNet with temporal convolutions, and VideoPhy2 \cite{bansal2025videophy2} relies on a 7B-parameter Transformer enhanced with LoRA adapters; all three operate with a 768‐dimensional hidden state, a shorter 77‐token prompt window, and varied positional embeddings (RoPE, Fourier, or sinusoidal).
This table highlights how model size, hidden dimension, position‐encoding strategy, and maximum prompt trade off across different approaches to physics-aware text-to-video synthesis and evaluation.

\begin{table*}[h]
\centering
\caption{Model details reported across the four papers.}
\label{tab:details}
% \small
\resizebox{\linewidth}{!}{
    \begin{tabular}{llllll}
    \toprule
    \textbf{Model} & \textbf{Core architecture} & \textbf{Size} & \textbf{Hidden Dim} & \textbf{Pos Enc} & \textbf{Max Prompt}\\
    \midrule
    CogVideoX-5B & 3D causal VAE $+$  Transformer & 5B & 4096 & RoPE & 226 \\
    CogVideoX-2B & 3D causal VAE $+$  Transformer & 2B & 4096 & sinusoidal & 226 \\
    Lavie & Cascaded Video Latent Diffusion  & 860M & 768 & RoPE & 77 \\
    VideoCrafter 2 &  UNet + temporal convolution & 1.2B & 768 & Fourier  & 77 \\
    VideoPhy2 & 7B Transformer with LoRA & 7B & 768 & sinusoidal & 77   \\
    \bottomrule
    \end{tabular}
}
\end{table*}
